\section{Introduction}
\label{sec:intro}

The proliferation of mobile devices and location-based services has
produced an abundance of spatial-temporal mobility data, which in turn
enables applications in urban planning~\cite{zheng2014urban,
jayarajah2018understanding}, transportation-equity analysis, contagious
disease management~\cite{eubank2004modelling,
wesolowski2012quantifying, ghadiri2024recreational}, and next-location
prediction more broadly~\cite{luca2021survey}. These same mobility
traces pose serious re-identification risk, since individuals are often
identifiable from their movement patterns
alone~\cite{de2013unique, saleheen2016msieve}. Privacy-utility trade-off
(PUT) models learn representations that preserve downstream utility
while suppressing re-identifiability, but prior work evaluating such
models has found that they can satisfy \emph{group} fairness checks
while violating \emph{individual} fairness -- treating users with
similar mobility behavior very differently in how much privacy
protection they actually receive~\cite{zhan2023analysing}. Equity
analysis in transportation research has traditionally relied on
coarser, demographic or region-level disparity
measures~\cite{yan2020fairness, yan2019fairst, he2025framework}; this
paper instead asks the question at the level of the individual user,
following the individual-fairness principle that similar individuals
should receive similar treatment~\cite{dwork2012fairness}.

A learned mobility-privacy mechanism is typically evaluated by training
it once, attacking its output with a single, group-agnostic
re-identification attacker, and reporting an aggregate or per-group
privacy gain. This paper asks a narrower question than ``is the
mechanism private'': \emph{who, specifically, is protected, and does
the mechanism's own evaluation protocol correctly report that}? If this
question is answered with an unreliable measurement protocol, a
mechanism can appear equitable while in fact systematically
under-protecting exactly the users most in need of protection.

We study this question on a representative learned mechanism from the
mobility-privacy literature, referred to throughout this paper in the
third person as Mo-PAE~\cite{zhan2023mopae} (an encoder/decoder/adversary
architecture that trains a feature encoder adversarially against an
in-model re-identification head, following the published PETS 2023
design this project audited and reimplemented in PyTorch).

\textbf{Contributions.} This paper makes four contributions:
\begin{itemize}
\item \textbf{A disparate-risk finding} (Section~\ref{sec:disparate-risk}):
  across three real mobility datasets and four Gowalla metro subsets,
  Mo-PAE's privacy gain is uneven across users in a specific, consistent
  direction -- the lowest-entropy decile, the most routine and so the
  most re-identifiable to begin with, receives the smallest relative
  risk reduction of any decile and is left with a residual
  re-identification risk 2--9$\times$ higher than the best-protected
  decile, a disparity that survives standard regression controls and is
  robust to bootstrap resampling.
\item \textbf{A test of whether an explicit fairness term fixes this}
  (Section~\ref{sec:phase4}): adding a fairness term to the training
  objective and raising its weight does not reliably close the gap.
  Depending on the formulation and dataset, raising the fairness weight
  either leaves the most-identifiable decile's risk statistically
  unchanged while further protecting an already-better-protected
  decile, or trades real utility for risk reduction without improving
  on the undefended mechanism. No formulation we test is shown to
  dominate plain Mo-PAE -- higher utility \emph{and} lower
  most-identifiable-decile risk, both beyond measurement noise -- once
  training is run to convergence and encoder-initialization variance is
  properly accounted for.
\item \textbf{A mechanistic account of why} (Section~\ref{sec:residual-risk}):
  a fairness term built on the mechanism's own co-trained adversary is
  targeting the wrong signal. That adversary's confidence stays close
  to chance and nearly flat across entropy deciles throughout training
  -- because the encoder is optimized specifically to suppress it --
  while a fresh external attacker evaluated on the same released
  features is both far more accurate and clearly decile-dependent. A
  fairness term that only ever sees the first signal cannot reliably
  target the risk the second one reveals.
\item \textbf{A robustness check on an adjacent, attacker-side claim}
  (Appendix~\ref{sec:adaptive-appendix}): we revisit the hypothesis that
  a fairness-aware \emph{adaptive} adversary -- one that knows which
  behavioral group a mechanism protects most and specializes its attack
  accordingly -- recovers substantially more risk than a naive,
  group-agnostic attacker. We find this is almost entirely an artifact
  of two confounds that require no attacker sophistication: a smaller
  per-group candidate pool, and oracle knowledge of a user's true
  behavioral group that a real adversary does not have. Once both are
  controlled for, a realistic adaptive adversary -- one that must infer
  a user's group from the same released features, achieving only
  34--42\% group-classification accuracy against a 10\% chance baseline
  -- performs no better, and in most cases worse, than the naive
  attacker it was meant to improve on.
\end{itemize}

Together, these results point toward a specific, falsifiable claim
about learned mobility-privacy mechanisms, not a general indictment of
the approach: individual-fairness claims about such a mechanism must be
measured at the level of who is actually protected using an attacker
that cannot be gamed by the mechanism's own training objective, and a
fairness term added to fix an uneven mechanism must be checked against
that same measurement before its effect is taken as improvement.
